\chapter{KẾT QUẢ THỰC NGHIỆM}
\label{chap:experiments}

\section{Thiết lập thử nghiệm}

\subsection{Tập dữ liệu}

Chúng tôi đánh giá mô hình trên các tập dữ liệu sau:

\begin{table}[H]
\centering
\caption{Thống kê tập dữ liệu}
\label{tab:datasets}
\begin{tabular}{lccc}
\toprule
Dataset & Train & Test & Classes \\
\midrule
Dataset A & 10.000 & 2.000 & 3 \\
Dataset B & 15.000 & 3.000 & 3 \\
\bottomrule
\end{tabular}
\end{table}

\subsection{Các chỉ số đánh giá}

Sử dụng các chỉ số sau để đánh giá hiệu suất mô hình:
\begin{itemize}
    \item Độ chính xác (Accuracy)
    \item Precision, Recall, F1-score
    \item Ma trận nhầm lẫn (Confusion Matrix)
\end{itemize}

\subsection{Siêu tham số}

Bảng~\ref{tab:hyperparams} cho thấy các siêu tham số được sử dụng trong thử nghiệm.

\begin{table}[H]
\centering
\caption{Cấu hình siêu tham số}
\label{tab:hyperparams}
\begin{tabular}{lc}
\toprule
Hyperparameter & Value \\
\midrule
Batch size & 32 \\
Learning rate & 0.001 \\
Epochs & 50 \\
Optimizer & Adam \\
\bottomrule
\end{tabular}
\end{table}

\section{Kết quả}

\subsection{Kết quả chính}

Bảng~\ref{tab:results} trình bày so sánh hiệu suất của các mô hình khác nhau.

\begin{table}[H]
\centering
\caption{So sánh hiệu suất}
\label{tab:results}
\begin{tabular}{lccc}
\toprule
Model & Accuracy & Precision & F1-score \\
\midrule
Baseline & 75,2 & 74,8 & 74,5 \\
Our Model & \textbf{82,5} & \textbf{81,9} & \textbf{82,1} \\
\bottomrule
\end{tabular}
\end{table}

\subsection{Ablation Study}

Chúng tôi thực hiện ablation study để hiểu đóng góp của từng thành phần.

\section{Thảo luận}

\subsection{Phân tích lỗi}

Phân tích các trường hợp lỗi phổ biến và xác định các khu vực cần cải thiện.
